% Generated draft: requires mathematical and source review.
\documentclass[11pt]{article}
\usepackage[margin=25mm]{geometry}
\usepackage{fontspec}
\setmainfont{Latin Modern Roman}
\usepackage{amsmath,amssymb,mathtools,unicode-math}
\setmathfont{Latin Modern Math}
\usepackage{enumitem,needspace,xurl,hyperref}
\setlength{\parindent}{0pt}
\setlength{\parskip}{5pt}
\setlength{\emergencystretch}{3em}
\newcommand{\N}{\mathbb{N}}
\newcommand{\Z}{\mathbb{Z}}
\newcommand{\Q}{\mathbb{Q}}
\newcommand{\R}{\mathbb{R}}
\newcommand{\C}{\mathbb{C}}
\newcommand{\F}{\mathbb{F}}
\newcommand{\A}{\mathbb{A}}
\newcommand{\PP}{\mathbb{P}}
\newcommand{\E}{\mathbb{E}}
\newcommand{\eps}{\varepsilon}
\newcommand{\dd}{\,\mathrm{d}}
\begin{document}
\section*{UnsolvedMath ID: 701774}
\textbf{Source problem code:} catalogue entry
\subsection{Definitions and mathematical statement}
A profinite group is a topological group isomorphic, as a topological group, to an inverse limit of finite discrete groups. Such a group \(G\) is torsion if every element has finite order, that is,
\[
\forall g\in G\;\exists n\in\Z_{\geq 1}\quad g^n=1_G.
\]
It has finite exponent if there is a single positive integer annihilating every element:
\[
\exists N\in\Z_{\geq 1}\;\forall g\in G\quad g^N=1_G.
\]
The integer \(N\) may depend on \(G\).

\textbf{Problem (catalogue entry).} Does every torsion profinite group have finite exponent?

The supplied source calls this a well-known problem and states that it reduces to the case of pro-\(p\) groups, citing W. Herfort. Here, for a prime \(p\), a pro-\(p\) group is a topological group isomorphic to an inverse limit of finite groups of \(p\)-power order. The reduced question asks, for every prime \(p\), whether every torsion pro-\(p\) group has finite exponent; the exponent bound may depend on the group and on \(p\). The reduction is reported from the supplied source without proof or independent verification. The source also refers to Archive 3.41, whose contents were not supplied.
\subsection{Short English statement}
Must the finite orders of elements in any torsion profinite group all divide one positive integer depending on the group? The supplied record labels this problem open and reports a reduction to pro-\(p\) groups; these source assessments have not been independently verified.
\subsection{Sources}
\begin{itemize}
\item Supplied local UnsolvedMath registry record, permanent ID \(701774\), exact source code \texttt{catalogue entry}, titled \emph{Kourovka Notebook Problem 20.116}. The record labels the problem open and attributes its proposal to John S. Wilson in 2022; the statement describes it as well-known.
\item \emph{The Kourovka Notebook (Unsolved Problems in Group Theory)}, Problem 20.116, 20th issue (2022), also identified in the supplied background with the 21st issue edition (2026), arXiv:1401.0300v48. The supplied collection description names E. I. Khukhro and V. D. Mazurov as editors.
\item W. Herfort, \emph{Arch. Math.}, \textbf{33} (1980), 404--410, cited by the supplied statement for the reduction to pro-\(p\) groups. No article title or further bibliographic details were supplied.
\item Archive 3.41, cross-reference supplied by the problem statement; the referenced text was not supplied.
\end{itemize}
\end{document}
